\bookchapter{Capítulo 16: Sistema de contas e serviços em nuvem}

\chapterintro{Um amigo meu teve o celular roubado no ano passado. A primeira coisa que ele pensou não foi “quanto custa o celular”, mas sim “ainda consigo recuperar minhas fotos?”. Infelizmente, ele não havia ativado o backup em nuvem. O celular é um bem material, mas os dados na sua conta são coisas que, uma vez perdidas, nunca mais voltam. Neste capítulo, vamos conversar sobre como fazer um verdadeiro seguro para sua vida digital.}

\booksection{I. Sua conta do celular: quantas coisas ela realmente gerencia?}

Muita gente acha que a conta serve só para baixar apps. Na verdade, uma conta de celular está vinculada a muito mais coisas do que você imagina:

\begin{itemize}
  \item \textbf{Contatos}: Números de telefone que você acumulou por anos, até décadas.
  \item \textbf{Fotos e vídeos}: Do nascimento dos filhos à formatura, do primeiro encontro à viagem de lua de mel.
  \item \textbf{Notas}: Inspirações que você anotou, diários, senhas de contas (sim, muita gente anota senhas nas Notas).
  \item \textbf{Histórico de compras de apps}: Centenas de reais gastos em apps pagos e compras dentro de apps.
  \item \textbf{Buscar celular}: Depois de perder o telefone, esta é a única ferramenta que pode ajudar a localizar, bloquear e apagar remotamente os dados.
\end{itemize}

Se você usa iPhone, tudo isso está vinculado ao seu Apple ID. Se usa Android, está vinculado à sua conta Huawei/Xiaomi/OPPO/vivo e também à sua conta Google. Essas contas são o “registro de identidade” da sua vida digital. Perder o registro de identidade deixa muitas coisas sem explicação.

\booksection{II. Três coisas obrigatórias: faça agora}

\begin{enumerate}
  \item \textbf{Ative a autenticação de dois fatores}: Seja você usuário de Apple ou Android, a primeira coisa a fazer é ir nas configurações da conta e ativar a "Autenticação de dois fatores" ou "Verificação em duas etapas". É como colocar uma fechadura extra na sua porta — só a senha não basta, é preciso que seu telefone receba um código de verificação. Mesmo que alguém adivinhe sua senha online, sem seu telefone, não conseguirá entrar.
  \item \textbf{Vincule um método de recuperação confiável}: A autenticação de dois fatores é muito segura, mas tem um risco — e se você perder seu próprio telefone? Por isso o sistema solicita que você configure uma "Chave de recuperação" ou "E-mail/telefone de backup". Esta é a única forma de recuperar sua conta em uma emergência. É melhor anotar em papel, não guardar no celular.
  \item \textbf{Confirme que o "Buscar celular" está ativado}: Esta função, depois que o celular é perdido, ajuda a localizar, bloquear, emitir som e até apagar remotamente todos os dados. É a última linha de defesa. Muita gente pula essa etapa ao comprar o celular, e quando precisa, descobre que não estava ativado — aí é tarde demais.
\end{enumerate}

\booksection{III. iCloud: o “ar” dos usuários Apple}

Quem usa iPhone, mais cedo ou mais tarde, enfrenta uma pergunta crucial: o que fazer quando o espaço do iCloud acaba?

A Apple oferece 5GB grátis. Sinceramente, esses 5GB não são suficientes nem para um backup. Por isso a maioria das pessoas acaba recebendo aquele pop-up irritante: “Armazenamento do iCloud Cheio”. Nesse momento, você tem três opções:

\begin{itemize}
  \item \textbf{Limpar}: Vá em Ajustes > Seu nome > iCloud > Gerenciar Armazenamento e veja quem está ocupando seu espaço. Geralmente, fotos e backups são os maiores consumidores. Comece excluindo backups antigos desnecessários.
  \item \textbf{Fazer upgrade}: 50GB por R\$ 3,50/mês, 200GB por R\$ 10,90/mês. Pessoalmente, acho que 50GB é suficiente para a maioria das pessoas, principalmente para backup e miniaturas de fotos.
  \item \textbf{Solução alternativa}: Se não quiser pagar, exporte as fotos periodicamente para o computador e desative a sincronização de fotos do iCloud. Mas, nesse caso, você perde a conveniência do backup automático de fotos.
\end{itemize}

Um lembrete importante: a Fototeca do iCloud não é um “drive em nuvem” — é uma “sincronização”. Se você excluir uma foto do celular, ela também será excluída do iCloud. Se o que você quer é “fazer backup na nuvem e depois liberar espaço no celular”, a Fototeca do iCloud não é o que você precisa. Você precisa do Backup do iCloud ou de um drive em nuvem de terceiros.

\begin{tcbinfo}{iCloud+: o trio de privacidade exclusivo para assinantes}
Basta fazer upgrade para qualquer plano pago do iCloud (50GB/200GB/2TB etc.) para se tornar automaticamente usuário\textbf{iCloud+}, desbloqueando três funções de proteção de privacidade:

\begin{itemize}
  \item \textbf{Private Relay (Retransmissão Privada)}: A "VPN" da Apple. O tráfego de navegação no Safari é retransmitido pelos servidores da Apple antes de chegar ao site. Ninguém (nem mesmo a Apple) consegue saber simultaneamente quem você é e o que você acessou. Ideal para proteger a privacidade em Wi-Fi público. Caminho:\texttt{Ajustes > Seu nome > iCloud > Retransmissão Privada}.
  \item \textbf{Hide My Email (Ocultar E-mail)}: Gera endereços de e-mail aleatórios descartáveis ao se registrar em sites/apps. Os e-mails são automaticamente encaminhados para seu e-mail real. Se não quiser mais receber e-mails daquele site, basta desativar o endereço aleatório, sem expor seu e-mail real. Caminho:\texttt{Ajustes > Seu nome > iCloud > Ocultar E-\allowbreak mail}.
  \item \textbf{Domínio de e-mail personalizado}: Você pode usar seu próprio domínio (como @seuSobrenome.com) para enviar e receber e-mails do iCloud, ideal para marca pessoal ou gerenciamento de e-mails familiares.
\end{itemize}

Nota: O Private Relay pode não estar disponível em alguns países/regiões. Usuários nesses locais podem precisar de soluções complementares.

\end{tcbinfo}

\begin{tcbwarn}{Proteção Avançada de Dados do iCloud (ADP) — criptografia de ponta a ponta de verdade}
No final de 2022, a Apple lançou a\textbf{Proteção Avançada de Dados do iCloud (Advanced Data Protection, ADP)}, expandindo a criptografia de ponta a ponta (E2EE) dos poucos dados originais (saúde, senhas) para\textbf{a grande maioria dos dados do iCloud}— fotos, notas, backup do iCloud, informações do dispositivo, totalizando 23 categorias de dados.

\begin{itemize}
  \item \textbf{O que é E2EE}: A chave de criptografia existe apenas nos seus dispositivos. Nem os servidores da Apple a possuem. Mesmo que a Apple seja hackeada ou que autoridades solicitem os dados, a Apple não consegue fornecer dados descriptografados legíveis.
  \item \textbf{Caminho para ativar}:\texttt{Ajustes > Seu nome > iCloud > Proteção Avançada de Dados}, siga as instruções para configurar contato de recuperação/chave de recuperação (\textbf{obrigatório}, se perder a chave, seus dados não poderão ser recuperados).
  \item \textbf{Observações}: (1)Após ativar, guarde a chave de recuperação com muito cuidado, caso contrário os dados não poderão ser recuperados em caso de perda da conta; (2)Alguns recursos como álbuns compartilhados e biblioteca compartilhada do iCloud não participam da E2EE; (3)Apple ID da China continental ainda não suporta ADP, sendo necessário trocar para conta de outra região.
\end{itemize}

Resumo em uma frase: Se sua conta contém dados privados importantes e a região da conta suporta, é altamente recomendável ativar o ADP — esta é a proteção de dados em nuvem mais forte que uma pessoa comum pode obter.

\end{tcbwarn}

\booksection{IV. Sincronização e backup: não confunda}

Essas duas palavras soam parecidas, mas as consequências são completamente diferentes:

\begin{itemize}
  \item \textbf{Sincronização}é bidirecional: você exclui no celular, exclui na nuvem; você exclui na nuvem, exclui no celular. Fotos do iCloud, Contatos, Notas — tudo isso é sincronização. A vantagem é a consistência entre dispositivos; a desvantagem é que "se excluiu, desapareceu mesmo".
  \item \textbf{Backup}é unidirecional: o celular tira um "instantâneo" do estado atual e o armazena na nuvem. Depois, se você excluir algo no celular, aquilo ainda estará no backup da nuvem. O Backup do iCloud e o backup local no computador são backups. A vantagem é que os dados ficam mais seguros; a desvantagem é que ocupa espaço e não é atualizado em tempo real.
\end{itemize}

\booksection{V. A matriz de contas dos usuários Android}

No lado Android, a situação é mais complexa que na Apple, porque seus dados podem estar dispersos em várias contas:

\begin{itemize}
  \item \textbf{Conta da marca}(Huawei/Xiaomi/OPPO/vivo): Gerencia seu espaço em nuvem, busca de celular, loja de aplicativos, carteira. Cada marca tem seu próprio serviço de nuvem, geralmente no modelo básico grátis + expansão paga.
  \item \textbf{Conta Google}: Se você usa Gmail, Google Fotos, Google Contatos, tudo isso está na conta Google. O Google oferece 15GB grátis, mais generoso que a Apple, mas é compartilhado com Gmail e Drive.
  \item \textbf{WeChat/QQ}: São independentes do sistema do celular. Seu histórico de conversas e favoritos do WeChat não estão na conta da marca nem na conta Google — estão apenas nos servidores do próprio WeChat.
\end{itemize}

Portanto, ao fazer backup no Android, é preciso verificar separadamente: backup em nuvem da marca, backup do Google e backup do histórico de conversas do WeChat. Nenhum dos três pode faltar.

\booksection{VI. Compartilhamento familiar: faça sua família também se beneficiar}

O “Compartilhamento Familiar” da Apple permite que você e sua família compartilhem compras de apps, músicas, filmes, espaço do iCloud e também a localização. Até 6 pessoas podem ser adicionadas. Configuração: Ajustes > Seu nome > Compartilhamento Familiar. Huawei e Xiaomi também têm funções semelhantes, que podem ser encontradas nas configurações da conta como “Compartilhamento familiar” ou “Nuvem familiar”.

Um cenário muito prático: configure o compartilhamento familiar no celular dos pais. Assim você pode ver a localização deles a qualquer momento (com o consentimento deles). Se um idoso se perder ou o telefone ficar incontactável, você terá uma localização.

\begin{tcbwarn}{As consequências de ter a conta roubada são mais graves do que você imagina}
Se um hacker obtiver sua senha da conta, ele pode: bloquear remotamente seu celular e pedir resgate, ver suas fotos e contatos, fazer compras com seus métodos de pagamento. Mais problemático ainda: se sua conta estiver vinculada a muitas coisas, recuperá-la é um pesadelo por si só. Então, autenticação de dois fatores: ative agora.

\end{tcbwarn}

\begin{tcbpractice}{Prática: faça um check-up completo na sua conta}
\begin{enumerate}
  \item Abra as configurações do celular, entre na página da conta e confirme que a autenticação de dois fatores está ativada.
  \item Verifique se o "Número de telefone confiável" e a "Chave de recuperação" estão configurados. Anote a chave em papel.
  \item Entre no iCloud ou gerenciamento de espaço em nuvem, veja quais dados estão sincronizando e desative os desnecessários.
  \item Confirme que o "Buscar meu celular" está ativado.
  \item Se você usa Android, verifique se o backup em nuvem da marca, backup do Google e backup do WeChat estão todos cobertos.
\end{enumerate}

\end{tcbpractice}
